\documentclass[10pt,a4paper]{article}
\usepackage{amsmath,amssymb,amsthm,graphicx}
\usepackage{hyperref}
\usepackage{fontspec}
\title{Explicit VP set-multilinear polynomial with near-optimal depth-sensitive formula lower bounds}
\author{Deepanshu Kush, Shubhangi Saraf}
\date{doi:10.4230/LIPIcs.CCC.2023.15; CC BY 4.0}
\begin{document}
\maketitle
\noindent\textit{Editorial scope.} Counted original Theorem10 for explicit F\_n,d: one construction with low-depth and arbitrary-depth lower bounds. Complete author Sections2–3 retained; stronger Section4 G\_n,d result is not claimed. Original field context and cited background lemmas are retained.\par\medskip
\noindent\textit{Allowed standard prerequisites (editorial).}\begin{itemize}
\item Symmetric-word relative-rank bounds of Kush–Saraf [22] and elementary rank algebra. Exact imported Lemmas11–12 and relative-rank properties are retained in the transcription; author does not reprove these cited background bounds.
\end{itemize}\medskip
% Faithful transcription of official CCC2023 article15, Sections2-3 (pp.15:11-15:15).
% Original author numbering is written explicitly; no new proofs or hypotheses.
\section*{2 Preliminaries}
\subsection*{2.1 Relative Rank and its Properties}
We first describe the notation that we need to define the measures that we use to prove Theorems 1 and 3.
\paragraph{Definition 5 (Relative Rank Measure of [23,36]).} Let $f$ be a polynomial that is set-multilinear with respect to the variable partition $(X_1,X_2,\ldots,X_d)$ where each set is of size $n$. Let $w=(w_1,w_2,\ldots,w_d)$ be a tuple (or word) of non-zero real numbers such that $2^{|w_i|}\in[n]$ for all $i\in[d]$. For each $i\in[d]$, let $X_i(w)$ be the variable set obtained by removing arbitrary variables from the set $X_i$ such that $|X_i(w)|=2^{|w_i|}$, and let $X(w)$ denote the tuple of sets of variables $(X_1(w),\ldots,X_d(w))$. Corresponding to a word $w$, define $P_w:=\{i\mid w_i>0\}$ and $N_w:=\{i\mid w_i<0\}$. Let $M_w^P$ be the set of all set-multilinear monomials over a subset of the variable sets $X_1(w),X_2(w),\ldots,X_d(w)$ indexed by $P_w$, and similarly let $M_w^N$ be the set of all set-multilinear monomials over these variable sets indexed by $N_w$.
Define the `partial derivative matrix' matrix $M_w(f)$ whose rows are indexed by the elements of $M_w^P$ and columns indexed by the elements of $M_w^N$ as follows: the entry of this matrix corresponding to a row $m_1$ and a column $m_2$ is the coefficient of the monomial $m_1\cdot m_2$ in $f$. We define
\[\operatorname{relrk}_w(f):=\frac{\operatorname{rank}(M_w(f))}{\sqrt{|M_w^P|\cdot|M_w^N|}}=\frac{\operatorname{rank}(M_w(f))}{2^{\frac12\sum_{i\in[d]}|w_i|}}.\]
\paragraph{Definition 6.} For any tuple $w=(w_1,\ldots,w_t)$ and a subset $S\subseteq[t]$, we shall refer to the sum $\sum_{i\in S}w_i$ by $w_S$. And by $w|_S$, we will refer to the tuple obtained by considering only the elements of $w$ that are indexed by $S$. We denote by $\mathbb F_{\rm sm}[T]$ the set of set-multilinear polynomials over the tuple of sets of variables $T$.
The following is a simple result that establishes various useful properties of the relative rank measure.
\paragraph{Claim 7 ([23]).}
\begin{enumerate}
\item (Imbalance) Say $f\in\mathbb F_{\rm sm}[X(w)]$. Then, $\operatorname{relrk}_w(f)\leq2^{-|w_{[d]}|/2}$.
\item (Sub-additivity) If $f,g\in\mathbb F_{\rm sm}[X(w)]$, then $\operatorname{relrk}_w(f+g)\leq\operatorname{relrk}_w(f)+\operatorname{relrk}_w(g)$.
\item (Multiplicativity) Say $f=f_1f_2\cdots f_t$ and assume that for each $i\in[t]$, $f_i\in\mathbb F_{\rm sm}[X(w|_{S_i})]$, where $(S_1,\ldots,S_t)$ is a partition of $[d]$. Then
\[\operatorname{relrk}_w(f)=\prod_{i\in[t]}\operatorname{relrk}_{w|_{S_i}}(f_i).\]
\end{enumerate}
\subsection*{2.2 Inner Product Gadget}
We crucially need the following observation to construct the hard polynomials in Theorems 1 and 3.
\paragraph{Observation 8.} Let $n=2^k$ and $X_1=\{x_{1,1},\ldots,x_{1,n}\}$ and $X_2=\{x_{2,1},\ldots,x_{2,n}\}$ be two disjoint sets of variables. Then, for any symmetric word $w\in\{k,-k\}^2$ (i.e., where $w_1+w_2=0$) and for the inner product ``gadget'' $f=X_1\cdot X_2=\sum_{i=1}^n x_{1,i}x_{2,i}$, $\operatorname{relrk}_w(f)=1$ i.e., $M_w(f)$ is full-rank.
\section*{3 A Hard Set-multilinear Polynomial in VP}
\subsection*{3.1 Description of the Polynomial}
Let $d$ be an even integer and let $X=(X_1,\ldots,X_d)$ be a collection of sets of variables where each $|X_i|=n$, and similarly, let $Y=(Y_1,\ldots,Y_d)$ be a distinct collection of sets of variables where each $|Y_i|=n$. We shall refer to the $Y$-variables as the auxiliary variables. For $i$ and $j\in\{1,\ldots,d\}$, let $X_i\cdot X_j$ denote the inner-product quadratic form $\sum_{k=1}^n x_{ik}x_{jk}$. Here, we shall assume that $X_i=\{x_{i,1},\ldots,x_{i,n}\}$ and $Y_i=\{y_{i,1},\ldots,y_{i,n}\}$.
For two integers $i\in\mathbb N$ and $j\in\mathbb N$, we denote $[i,j]=\{k\in\mathbb N:i\leq k\text{ and }k\leq j\}$ and call such a set an interval. For every interval $[i,j]\subseteq[d]$, we define a polynomial $f_{i,j}(X,Y)\in\mathbb F_{\rm sm}[X_i,\ldots,X_j,Y_i,\ldots,Y_j]$ as follows:
\[f_{i,j}=\begin{cases}
y_{i,j}y_{j,i}(X_i\cdot X_j)&\text{if }j=i+1,\\
0&\text{if }j-i\text{ is even},\\
y_{i,j}y_{j,i}(X_i\cdot X_j)\cdot f_{i+1,j-1}+\displaystyle\sum_{r=i+1}^{j-1} f_{i,r}f_{r+1,j}&\text{otherwise}.
\end{cases}\]
\paragraph{Remark 9.} As described in Section 1.5, other than the use of the inner product gadget, one key difference between $f_{i,j}$ and the construction in [29] is that it uses fewer auxiliary variables. More specifically, while [29] had a ``fresh'' auxiliary variable for every choice of $i,r,j$ in the sum, we are unable to afford that not only because it destroys the set-multilinearity of the polynomial but most importantly, because of the aforementioned degree blow-up. This is also the reason why more straightforward attempts to ``set-multilinearize'' [29] such as by adding two ``copies'' $y_0$ and $y_1$ for each of their auxiliary variables $y$ (where intuitively $y_0$ and $y_1$ correspond to setting $y$ as $0$ or $1$ respectively in their argument) do not work.
The following is a more precise and general version of Theorem 3 that is stated in Section 1. We also incorporate Remark 4 here and show our lower bound for any degree $d\leq n$. Theorem 3 follows from the special case $d=n$.

\paragraph{Theorem 10.} Let $n=2^k$, and suppose $d\leq n$ be an even integer that is large enough\footnote{We only need $d$ to be larger than some absolute constant.}, and $1\leq\Delta\leq\log d/\log\log d$ be any positive integer. Let $X_i,Y_i$ denote the sets of $n$ variables $\{x_{i,j}:j\in[n]\}$ and $\{y_{i,j}:j\in[n]\}$ respectively and let $X,Y$ be the tuples $(X_1,\ldots,X_d)$ and $(Y_1,\ldots,Y_d)$. Then,
\begin{itemize}
\item there is a $\operatorname{poly}(n,d)$-size set-multilinear circuit computing $F_{n,d}=f_{1,d}(X,Y)$ as defined above,
\item any set-multilinear formula of product-depth $\Delta$ computing $F_{n,d}$ must have size at least $d^{\Omega(d^{1/\Delta}/\Delta)}$, and
\item further, any set-multilinear formula of arbitrary product-depth computing $F_{n,d}$ must have size at least $d^{\Omega(\log d)}$.
\end{itemize}

\subsection*{3.2 Proof of Hardness}
Note that the first item in Theorem 10 follows immediately from the recursive definition of $f_{1,d}$ (notice that there are only up to $d^2$ many distinct intervals of $[d]$). For proving the next two items, we invoke the symmetric word framework of [22]. The following couple of lemmas help establish that the relative rank measure with respect to symmetric words is (suitably) small for low-depth and general-depth set-multilinear formulas, respectively.
\paragraph{Lemma 11 ([22]).} Let $C$ be a set-multilinear formula of product-depth $1\leq\Delta\leq\log d/\log\log d$ of size at most $s$ which computes a polynomial (over any fixed field) that is set-multilinear with respect to the partition $(X_1,\ldots,X_d)$ where each $|X_i|=n$. Let $w\in\{k,-k\}^d$ be chosen uniformly at random. Then, we have
\[\operatorname{relrk}_w(C)\leq s\cdot2^{-kd^{1/\Delta}/20}\]
with probability at least $1-s\cdot d^{-d^{1/\Delta}/(12\Delta)}$.
\paragraph{Lemma 12 ([22]).} Let $F$ be a set-multilinear formula of size at most $s$ which computes a polynomial (over any fixed field) that is set-multilinear with respect to the partition $(X_1,\ldots,X_d)$ where each $|X_i|=n$. Let $w\in\{k,-k\}^d$ be chosen uniformly at random. Then, we have
\[\operatorname{relrk}_w(F)\leq s\cdot2^{-k\log d/20}\]
with probability at least $1-s\cdot d^{-\log d/60}$.
Next, we shall show that the hard polynomial $F_{n,d}$ in Theorem 10 has high relative rank (in fact, the maximum possible value -- $1$) with respect to a symmetric word. For this, we consider an alternate view of these polynomials and require the following notion. For an even integer $d$, define $\operatorname{Dyck}(d)$ to be the collection of all strings (called Dyck words) of length $d$ over symbols `(' and `)' that are well-matched in the natural way. More precisely, it is the collection of all strings $u$ of length $d$ such that all prefixes of $u$ contain no more $)$'s than $($'s and the total number of $($'s in $u$ equals the total number of $)$. For example, ``()()'' and ``(())'' belong to $\operatorname{Dyck}(4)$ but not ``(()(''. Note that for any `(' appearing in a Dyck word, there is a unique `)' which ``closes'' it. Given a Dyck word $u\in\operatorname{Dyck}(d)$, we call $(i,j)$ a matching parenthesis pair of $u$ if there is `(' in the $i$-th position of $u$ that is closed by a `)' in the $j$-th position of $u$ (clearly then, $j-i>0$ must be odd).
Given a string $u\in\operatorname{Dyck}(d)$ and the setup above for defining the polynomials $f_{i,j}$, we can associate to $u$ a product of inner products polynomial $\operatorname{IP}_u\in\mathbb F_{\rm sm}[X_i,\ldots,X_j]$ in the natural way: define $\operatorname{IP}_u$ to be the product of all $X_i\cdot X_j$ where $(i,j)$ is a matching parenthesis pair of $u$. For example, the strings ``()()'' and ``(())'' would correspond to the polynomials $(X_1\cdot X_2)\cdot(X_3\cdot X_4)$ and $(X_1\cdot X_4)\cdot(X_2\cdot X_3)$ respectively. We define $y_u$ analogously: it is the product of all $y_{i,j}\cdot y_{j,i}$ where $(i,j)$ is a matching parenthesis pair of $u$. So, if $u=\text{``(())''}\in\operatorname{Dyck}(4)$, then $y_u=y_{1,4}y_{4,1}y_{2,3}y_{3,2}$. The following observation then follows immediately from the recursive definition of $f_{i,j}$.
\paragraph{Observation 13.} For every interval $[i,j]\subseteq[d]$ where $j-i$ is odd, there exist constants $c_u\in\mathbb F$ corresponding to every $u\in\operatorname{Dyck}(j-i+1)$ such that
\[f_{i,j}(X,Y)=\sum_{u\in\operatorname{Dyck}(j-i+1)}c_u y_u\operatorname{IP}_u.\]\footnote{Strictly speaking, the indices within $\operatorname{IP}_u$ and $y_u$ here need to be ``translated'' appropriately to suit the interval $[i,j]$ (which may not necessarily be $[1,j-i+1]$).}
Moreover, if $\mathbb F$ has characteristic zero, then every $c_u\neq0$.
\paragraph{Claim 14.} Let $d$ be a positive even integer. For any $w\in\{-k,k\}^d$ with $w_{[d]}=0$ (i.e., $|P_w|=|N_w|$), there exists a Dyck work $u\in\operatorname{Dyck}(d)$ such that for every matching parenthesis pair $(i,j)$ of $u$, either $i\in P_w$ and $j\in N_w$, or $i\in N_w$ and $j\in P_w$.
\paragraph{Proof.} We prove this by induction on $d$. The base case $d=2$ is trivial as there is only a single matching parenthesis pair $(1,2)$ for which the given condition must indeed hold. Now, suppose $d>2$ and $w\in\{-k,k\}^d$ is a given word with $w_{[d]}=0$. Let us refer to $P_w$ and $N_w$ as the two ``parts'' of $w$ and take cases on the membership of $1$ and $d$ in these sets.
\begin{description}
\item[Case 1:] $1$ and $d$ are in different parts. By the induction hypothesis, there is a Dyck word $u'\in\operatorname{Dyck}(d-2)$ (corresponding to the subset $\{2,\ldots,d-1\}$ of indices) for which the desired condition holds. Hence, we can simply define $u$ to be the string ``$(u')$'' and the claim follows.
\item[Case 2:] $1$ and $d$ are in the same part. Notice that there exists an index $r$ such that both $w_{[1,r]}$ and $w_{[r+1,d]}$ are $0$. Then, we can define $u$ to be the concatenation of the two Dyck words that the induction hypothesis yields for the intervals $[1,r]$ and $[r+1,d]$ respectively and the claim follows.
\end{description}\hfill$\triangleleft$
\paragraph{Lemma 15.} Let $n=2^k$ and $d\leq n$ be an even integer. Over any field $\mathbb F$ of characteristic zero, the polynomial $F_{n,d}=f_{1,d}\in\mathbb F_{\rm sm}[X,Y]$ as defined above satisfies the following: For any $w\in\{-k,k\}^d$ with $w_{[d]}=0$, $M_w(F_{n,d})$ is full-rank when viewed as a matrix over the field $\mathbb F(Y)$, the field of rational functions over the $Y$ variables.
\paragraph{Proof.} Fix a word $w\in\{-k,k\}^d$ with $w_{[d]}=0$ and let $s\in\operatorname{Dyck}(d)$ be a Dyck word as given by Claim 14. By Observation 13, we know that $F=f_{1,d}$ has the form $\sum_{u\in\operatorname{Dyck}(d)}c_u y_u\operatorname{IP}_u$. Consider the polynomial $f$ obtained by plugging in $y_{i,j}=y_{j,i}=0$ for every $i,j$ such that $(i,j)$ is not a matching parenthesis pair of $s$, and $y_{i,j}=y_{j,i}=1$ for every $i,j$ such that $(i,j)$ is a matching parenthesis pair of $s$. Observe that the only surviving term from $F$ in $f$ is the one indexed by $s$. Therefore, to argue that $M_w(F)$ is full-rank over $\mathbb F(Y)$\footnote{We need to show that its determinant -- a polynomial in $\mathbb F[Y]$ -- is non-zero.}, it suffices to show that $M_w(c_s\operatorname{IP}_s)$ is full-rank. As $c_s\neq0$ by Observation 13, this follows from Observation 8 (the matrix of the inner product gadget has full rank), Claim 14 ($w$ ``splits'' every matching parenthesis pair of $s$), and Claim 7 (the multiplicativity of $\operatorname{relrk}_w$).\hfill$\triangleleft$
Let us return to the proof of the last two items of Theorem 10. Let $C$ be a set-multilinear formula of product depth $\Delta$ of size $s$ computing $F_{n,d}(X)$ (now interpreted as a formula over the field $\mathbb F(Y)$). Suppose $s<d^{d^{1/\Delta}/(24\Delta)}$. Then, by Lemma 11, with probability at least $1-d^{-d^{1/\Delta}/(24\Delta)}$,
\[\operatorname{relrk}_w(C)\leq s\cdot2^{-kd^{1/\Delta}/20}.\]
But now, we can condition on the event that $w_{[d]}=0$ (which occurs with probability $\Theta(1/\sqrt d)$) to establish the existence of a word $w\in\{-k,k\}^d$ with $w_{[d]}=0$ such that $w$ satisfies $\operatorname{relrk}_w(C)\leq s\cdot2^{-kd^{1/\Delta}/20}$. This is because of the asymptotic bound $1/\sqrt d\gg d^{-d^{1/\Delta}/(24\Delta)}$, which follows from the given constraints on the parameters $d,\Delta$. Therefore, by Lemma 15,
\[s\geq2^{kd^{1/\Delta}/20}\cdot\operatorname{relrk}_w(C)=n^{d^{1/\Delta}/20}\]
which contradicts the assumption that $s<d^{d^{1/\Delta}/(24\Delta)}$. Thus, we conclude that $s\geq d^{d^{1/\Delta}/(24\Delta)}=d^{\Omega(d^{1/\Delta}/\Delta)}$.
Similarly, to see the final item of Theorem 10, let $F$ be a set-multilinear formula of size $s$ computing $F_{n,d}$ (now interpreted as a formula over the field $\mathbb F(Y)$). Suppose $s<d^{\log d/120}$. Then, by Lemma 12, with probability at least $1-d^{-\log d/120}$,
\[\operatorname{relrk}_w(F)\leq s\cdot2^{-k\log d/20}.\]
But now, we can condition on the event that $w_{[d]}=0$ (which occurs with probability $\Theta(1/\sqrt d)$) to establish the existence of a word $w\in\{-k,k\}^d$ with $w_{[d]}=0$ such that $w$ satisfies $\operatorname{relrk}_w(F)\leq s\cdot2^{-k\log d/20}$. This is because of the trivial asymptotic bound $1/\sqrt d\gg d^{-\log d/120}$. Therefore, again by Lemma 15,
\[s\geq2^{k\log d/20}\cdot\operatorname{relrk}_w(F)=n^{\log d/20}\]
which contradicts the assumption that $s<d^{\log d/120}$. Thus, we conclude that $s\geq d^{\log d/120}=d^{\Omega(\log d)}$.

\begin{thebibliography}{38}
\bibitem[1]{ref1} Manindra Agrawal, Rohit Gurjar, Arpita Korwar, and Nitin Saxena. Hitting-sets for ROABP and sum of set-multilinear circuits. SIAM J. Comput., 44(3):669–697, 2015. doi:10.1137/140975103.
\bibitem[2]{ref2} Manindra Agrawal and V. Vinay. Arithmetic circuits: A chasm at depth four. In 49th Annual IEEE Symposium on Foundations of Computer Science, FOCS 2008, October 25-28, 2008, Philadelphia, PA, USA, pages 67–75. IEEE Computer Society, 2008. doi:10.1109/FOCS.2008.32. P1
\bibitem[3]{ref3} Miklós Ajtai. $\Sigma^1\_1$-formulae on finite structures. Ann. Pure Appl. Log., 24(1):1–48, 1983. doi:10.1016/0168-0072(83)90038-6.
\bibitem[4]{ref4} Vikraman Arvind and S. Raja. Some lower bound results for set-multilinear arithmetic computations. Chic. J. Theor. Comput. Sci., 2016, 2016. URL: http://cjtcs.cs.uchicago.edu/articles/2016/6/contents.html.
\bibitem[5]{ref5} Peter Bürgisser. Cook’s versus valiant’s hypothesis. Theor. Comput. Sci., 235(1):71–88, 2000. doi:10.1016/S0304-3975(99)00183-8.
\bibitem[6]{ref6} Suryajith Chillara, Christian Engels, Nutan Limaye, and Srikanth Srinivasan. A near-optimal depth-hierarchy theorem for small-depth multilinear circuits. In Mikkel Thorup, editor, 59th IEEE Annual Symposium on Foundations of Computer Science, FOCS 2018, Paris, France, October 7-9, 2018, pages 934–945. IEEE Computer Society, 2018. doi:10.1109/FOCS.2018.00092.
\bibitem[7]{ref7} Suryajith Chillara, Nutan Limaye, and Srikanth Srinivasan. Small-depth multilinear formula lower bounds for iterated matrix multiplication with applications. SIAM J. Comput., 48(1):70– 92, 2019. doi:10.1137/18M1191567.
\bibitem[8]{ref8} Zeev Dvir, Guillaume Malod, Sylvain Perifel, and Amir Yehudayoff. Separating multilinear branching programs and formulas. In Howard J. Karloff and Toniann Pitassi, editors, Proceed- ings of the 44th Symposium on Theory of Computing Conference, STOC 2012, New York, NY, USA, May 19–22, 2012, pages 615–624. ACM, 2012. doi:10.1145/2213977.2214034.
\bibitem[9]{ref9} Michael A. Forbes and Amir Shpilka. Quasipolynomial-time identity testing of non- commutative and read-once oblivious algebraic branching programs. In 54th Annual IEEE Symposium on Foundations of Computer Science, FOCS 2013, 26-29 October, 2013, Berkeley, CA, USA, pages 243–252. IEEE Computer Society, 2013. doi:10.1109/FOCS.2013.34.
\bibitem[10]{ref10} Hervé Fournier, Nutan Limaye, Guillaume Malod, and Srikanth Srinivasan. Lower bounds for depth-4 formulas computing iterated matrix multiplication. SIAM J. Comput., 44(5):1173–1201, 2015. doi:10.1137/140990280.
\bibitem[11]{ref11} Merrick L. Furst, James B. Saxe, and Michael Sipser. Parity, circuits, and the polynomial-time hierarchy. Math. Syst. Theory, 17(1):13–27, 1984. doi:10.1007/BF01744431.
\bibitem[12]{ref12} Zeyu Guo and Rohit Gurjar. Improved explicit hitting-sets for roabps. In Jaroslaw Byrka and Raghu Meka, editors, Approximation, Randomization, and Combinatorial Optimization. Al- gorithms and Techniques, APPROX/RANDOM 2020, August 17-19, 2020, Virtual Conference, volume 176 of LIPIcs, pages 4:1–4:16. Schloss Dagstuhl – Leibniz-Zentrum für Informatik, 2020. doi:10.4230/LIPIcs.APPROX/RANDOM.2020.4.
\bibitem[13]{ref13} Johan Håstad. Almost optimal lower bounds for small depth circuits. In Juris Hartmanis, editor, Proceedings of the 18th Annual ACM Symposium on Theory of Computing, May 28-30, 1986, Berkeley, California, USA, pages 6–20. ACM, 1986. doi:10.1145/12130.12132.
\bibitem[14]{ref14} Pavel Hrubes and Amir Yehudayoff. Homogeneous formulas and symmetric polynomials. Comput. Complex., 20(3):559–578, 2011. doi:10.1007/s00037-011-0007-3.
\bibitem[15]{ref15} Neeraj Kayal, Nutan Limaye, Chandan Saha, and Srikanth Srinivasan. An exponential lower bound for homogeneous depth four arithmetic formulas. SIAM J. Comput., 46(1):307–335, 2017. doi:10.1137/151002423.
\bibitem[16]{ref16} Neeraj Kayal, Vineet Nair, and Chandan Saha. Separation between read-once oblivious algebraic branching programs (roabps) and multilinear depth-three circuits. ACM Trans. Comput. Theory, 12(1):2:1–2:27, 2020. doi:10.1145/3369928.
\bibitem[17]{ref17} Neeraj Kayal, Chandan Saha, and Ramprasad Saptharishi. A super-polynomial lower bound for regular arithmetic formulas. In David B. Shmoys, editor, Symposium on Theory of Computing, STOC 2014, New York, NY, USA, May 31 – June 03, 2014, pages 146–153. ACM, 2014. doi:10.1145/2591796.2591847.
\bibitem[18]{ref18} Neeraj Kayal, Chandan Saha, and Sébastien Tavenas. An almost cubic lower bound for depth three arithmetic circuits. In Ioannis Chatzigiannakis, Michael Mitzenmacher, Yuval Rabani, and Davide Sangiorgi, editors, 43rd International Colloquium on Automata, Languages, and Programming, ICALP 2016, July 11-15, 2016, Rome, Italy, volume 55 of LIPIcs, pages 33:1–33:15. Schloss Dagstuhl – Leibniz-Zentrum für Informatik, 2016. doi:10.4230/LIPIcs.ICALP.2016.33.
\bibitem[19]{ref19} Neeraj Kayal, Chandan Saha, and Sébastien Tavenas. On the size of homogeneous and of depth-four formulas with low individual degree. Theory Comput., 14(1):1–46, 2018. doi:10.4086/toc.2018.v014a016.
\bibitem[20]{ref20} Pascal Koiran. Arithmetic circuits: The chasm at depth four gets wider. Theor. Comput. Sci., 448:56–65, 2012. doi:10.1016/j.tcs.2012.03.041.
\bibitem[21]{ref21} Mrinal Kumar and Shubhangi Saraf. On the power of homogeneous depth 4 arithmetic circuits. SIAM J. Comput., 46(1):336–387, 2017. doi:10.1137/140999335.
\bibitem[22]{ref22} Deepanshu Kush and Shubhangi Saraf. Improved low-depth set-multilinear circuit lower bounds. In Shachar Lovett, editor, 37th Computational Complexity Conference, CCC 2022, July 20-23, 2022, Philadelphia, PA, USA, volume 234 of LIPIcs, pages 38:1–38:16. Schloss Dagstuhl – Leibniz-Zentrum für Informatik, 2022. doi:10.4230/LIPIcs.CCC.2022.38.
\bibitem[23]{ref23} Nutan Limaye, Srikanth Srinivasan, and Sébastien Tavenas. Superpolynomial lower bounds against low-depth algebraic circuits. In 62nd IEEE Annual Symposium on Foundations of Computer Science, FOCS 2021, Denver, CO, USA, February 7-10, 2022, pages 804–814. IEEE, 2021. doi:10.1109/FOCS52979.2021.00083.
\bibitem[24]{ref24} Noam Nisan and Avi Wigderson. Hardness vs randomness. J. Comput. Syst. Sci., 49(2):149–167, 1994. doi:10.1016/S0022-0000(05)80043-1.
\bibitem[25]{ref25} Noam Nisan and Avi Wigderson. Lower bounds on arithmetic circuits via partial derivatives. Comput. Complex., 6(3):217–234, 1997. doi:10.1007/BF01294256.
\bibitem[26]{ref26} Ran Raz. Separation of multilinear circuit and formula size. Theory Comput., 2(6):121–135, 2006. doi:10.4086/toc.2006.v002a006.
\bibitem[27]{ref27} Ran Raz. Multi-linear formulas for permanent and determinant are of super-polynomial size. J. ACM, 56(2):8:1–8:17, 2009. doi:10.1145/1502793.1502797.
\bibitem[28]{ref28} Ran Raz. Tensor-rank and lower bounds for arithmetic formulas. J. ACM, 60(6):40:1–40:15, 2013. doi:10.1145/2535928.
\bibitem[29]{ref29} Ran Raz and Amir Yehudayoff. Balancing syntactically multilinear arithmetic circuits. Comput. Complex., 17(4):515–535, 2008. doi:10.1007/s00037-008-0254-0.
\bibitem[30]{ref30} Ran Raz and Amir Yehudayoff. Lower bounds and separations for constant depth multilinear circuits. Comput. Complex., 18(2):171–207, 2009. doi:10.1007/s00037-009-0270-8.
\bibitem[31]{ref31} Alexander A. Razborov. Lower bounds on the size of bounded depth circuits over a complete basis with logical addition. Mathematical notes of the Academy of Sciences of the USSR, 41:333–338, 1987.
\bibitem[32]{ref32} Ramprasad Saptharishi. A survey of lower bounds in arithmetic circuit complexity. Github Survey, 2015. URL: https://github.com/dasarpmar/lowerbounds-survey.
\bibitem[33]{ref33} Amir Shpilka and Amir Yehudayoff. Arithmetic circuits: A survey of recent results and open questions. Found. Trends Theor. Comput. Sci., 5(3-4):207–388, 2010. doi:10.1561/0400000039.
\bibitem[34]{ref34} Roman Smolensky. Algebraic methods in the theory of lower bounds for boolean circuit complexity. In Alfred V. Aho, editor, Proceedings of the 19th Annual ACM Symposium on Theory of Computing, 1987, New York, New York, USA, pages 77–82. ACM, 1987. doi:10.1145/28395.28404.
\bibitem[35]{ref35} Sébastien Tavenas. Improved bounds for reduction to depth 4 and depth 3. Inf. Comput., 240:2–11, 2015. doi:10.1016/j.ic.2014.09.004.
\bibitem[36]{ref36} Sébastien Tavenas, Nutan Limaye, and Srikanth Srinivasan. Set-multilinear and non- commutative formula lower bounds for iterated matrix multiplication. In Stefano Le- onardi and Anupam Gupta, editors, STOC ’22: 54th Annual ACM SIGACT Symposium on Theory of Computing, Rome, Italy, June 20–24, 2022, pages 416–425. ACM, 2022. doi:10.1145/3519935.3520044.
\bibitem[37]{ref37} Leslie G. Valiant, Sven Skyum, Stuart J. Berkowitz, and Charles Rackoff. Fast parallel computation of polynomials using few processors. SIAM J. Comput., 12:641–644, 1983.
\bibitem[38]{ref38} Andrew Chi-Chih Yao. Separating the polynomial-time hierarchy by oracles (preliminary version). In 26th Annual Symposium on Foundations of Computer Science, Portland, Oregon, USA, 21-23 October 1985, pages 1–10. IEEE Computer Society, 1985. doi:10.1109/SFCS.1985.49.
\end{thebibliography}
\end{document}
